\documentclass[12pt]{article}
\usepackage[utf8]{inputenc}
\usepackage[T1]{fontenc}
\usepackage{hyperref}
\title{Rana Tharu Language Resources for K-12 Curriculum Development}
\author{Compiled by Claude Sonnet 5 (claude-sonnet-5) for Andr\'as Kornai}
\date{}
\begin{document}
\maketitle

\section{Introduction}
Rana Tharu (ISO 639-3 thr) is an Indo-Aryan language spoken in far-western Nepal (Kailali and Kanchanpur) and in India, with about 490,000 speakers. This report is based on a small number of searches. Linguistic documentation exists, but I found no graded textbook.

\section{Primary Educational Resources}
\begin{itemize}
\item World Vision Nepal developed a series of 15 Tharu-language early-grade readers for Rana Tharu communities. \url{https://www.wvi.org/stories/nepal/embracing-teaching-learning-materials-tharu-language}
\item Wikipedia states that mother-tongue education in Rana Tharu is provided up to grade three in Kanchanpur. I did not verify this from a primary source. \url{https://en.wikipedia.org/wiki/Rana_Tharu_language}
\end{itemize}

\section{Linguistic and Reference Materials}
\begin{itemize}
\item Dubi Nanda Dhakal, trilingual Rana Tharu--Nepali--English dictionary (Lincom, 2015, 143 pp.). It is described as the first dictionary of the language. Data came mainly from two Kanchanpur speakers. \url{https://linguistlist.org/issues/26/4679/}
\item LINGUIST List notice for a related Dhakal volume. \url{https://linguistlist.org/issues/26/650/}
\item Dhakal, ``Notes on Rana Tharu grammar''. \url{https://www.academia.edu/43908543}
\item Pitambar Bhatta, ``Rana Tharu Verb Morphology'', MA thesis, Tribhuvan University, 2018. \url{https://elibrary.tucl.edu.np/handle/123456789/21978}
\end{itemize}

\section{Religious and Community Resources}
None were found in the search results. A general article on the community is at \url{https://en.wikipedia.org/wiki/Rana_Tharu_people}.

\section{Digital and Audio Resources}
None were found in the search results.

\section{Curriculum Development Context}
The dictionary and grammar notes give a base for orthography and vocabulary work. Early-grade readers exist and could be adapted. Textbooks beyond the early grades, and the current Ministry of Education position, need checking with the Curriculum Development Centre and local Tharu organisations.
\end{document}